\documentclass{article}
\usepackage{amsmath}
\usepackage{xcolor}
\begin{document}

\textbf{Cap 02 - E.D. Primer Orden}

Representacion de una E.D.


\begin{gather*}
f(x,y) = \frac{dy}{dx}   \rightarrow \text{ Forma General }(a) \\
sen(x)e^{x}+y = \frac{dy}{dx} \\
M(x,y)dx +N(x,y)dy = 0  \rightarrow \text{ Forma Diferencial }(b) \\
(e^{2x}y)dx+(x+y)dy = 0  \\
\vspace{0.5em} \\
\text{es posible llevar de una forma }(a)\text{  la forma }(b)\text{ y viceversa} \\
(e^{2x}y)dx+(x+y)dy = 0  \\
(e^{2x}y)+(x+y)\frac{dy}{dx} = 0  \\
(x+y)\frac{dy}{dx} = -e^{2x}y \\
\frac{dy}{dx} = \frac{-e^{2x}y}{(x+y)}
\end{gather*}


Tipos or Rutas de Solucion de acuerdo a su presentacion

\textbf{E.D. de Variables Separable}


\begin{gather*}
\text{\colorbox[RGB]{248,231,28}{$M$}}\text{\colorbox[RGB]{248,231,28}{$($}}\text{\colorbox[RGB]{248,231,28}{$x,y$}}\text{\colorbox[RGB]{248,231,28}{$)$}}\text{\colorbox[RGB]{248,231,28}{$dx +N$}}\text{\colorbox[RGB]{248,231,28}{$($}}\text{\colorbox[RGB]{248,231,28}{$x,y$}}\text{\colorbox[RGB]{248,231,28}{$)$}}\text{\colorbox[RGB]{248,231,28}{$dy = 0 $}} \\
\text{Si se puede garantizar que:} \\
M(x,y)  \rightarrow \text{ solo depende de x } \rightarrow  A(x) \\
y \\
N(x,y) \rightarrow \text{ solo depende de y } \rightarrow  B(y): \\
A(x)dx +B(y)dy = 0  \\
\text{la E.D es de variable separable, por ende se puede resolver integrando} \\
\int A(x)dx +\int B(y)dy = c  \rightarrow \text{ donde c  es una ctte } \\
\frac{dy}{dx} = \frac{x^{2}+2}{e^{y}}  \rightarrow   \\
e^{-y}dy = (x^{2}+2)dx \\
(x^{2}+2)dx-e^{-y}dy =0 \\
\text{Multiplica, division de factores uniamente de una variable individual}
\end{gather*}


\textbf{E.D Homogeneas $\rightarrow$  }Comunmente representadas en su forma (b)


\begin{gather*}
M(x,y)dx +N(x,y)dy = 0  \\
\text{Una forma es utiliza el operando }\text{\colorbox[RGB]{248,231,28}{$𝜆$}} \\
M(\text{\colorbox[RGB]{248,231,28}{$𝜆$}}x,\text{\colorbox[RGB]{248,231,28}{$𝜆$}}y)dx +N(\text{\colorbox[RGB]{248,231,28}{$𝜆$}}x,\text{\colorbox[RGB]{248,231,28}{$𝜆$}}y)dy = 0  \\
\text{Si es posible factorizar y lograr que }\text{\colorbox[RGB]{248,231,28}{$\text{𝜆 sea igual en ambos terminos}$}} \\
\text{la E.D. es homogenea}
\end{gather*}





\begin{gather*}
f(x,y) = \frac{dy}{dx}   \rightarrow \text{ Forma General }(a) \\
\frac{dy}{dx} = f(\text{\colorbox[RGB]{248,231,28}{$𝜆$}}x,\text{\colorbox[RGB]{248,231,28}{$𝜆$}}y)  \rightarrow f(x,y)  \rightarrow \text{ es Homogenea}
\end{gather*}


Ejemplo


\begin{gather*}
(x^{2}y+y^{3})dx+(x+y)dy = 0 \\
((x\text{\colorbox[RGB]{248,231,28}{$𝜆$}})^{2}y\text{\colorbox[RGB]{248,231,28}{$𝜆$}}+(y\text{\colorbox[RGB]{248,231,28}{$𝜆$}})^{3})dx+(\text{\colorbox[RGB]{248,231,28}{$𝜆$}}x+\text{\colorbox[RGB]{248,231,28}{$𝜆$}}y)dy = 0 \\
(x^{2}\text{\colorbox[RGB]{248,231,28}{$𝜆$}}\text{\colorbox[RGB]{248,231,28}{$^{3}$}}y+y^{3}\text{\colorbox[RGB]{248,231,28}{$𝜆$}}\text{\colorbox[RGB]{248,231,28}{$^{3}$}})dx+(𝜆x+𝜆y)dy = 0 \\
𝜆^{3}(x^{2}y+y^{3})dx+𝜆(x+y)dy = 0  \\
 \rightarrow \text{ no es Homogenea} \\
\vspace{0.5em} \\
(x^{2}y+y^{3})dx+(x^{3}+y^{2}x)dy = 0 \\
(x^{2}𝜆^{2}y𝜆+y^{3}𝜆^{3})dx+(x^{3}𝜆^{3}+y^{2}𝜆^{2}x𝜆)dy = 0 \\
𝜆^{3}(x^{2}y+y^{3})dx+𝜆^{3}(x^{3}+y^{2}x)dy = 0 \\
𝜆^{3}((x^{2}y+y^{3})dx+(x^{3}+y^{2}x)dy )=0 \\
 \rightarrow \text{ E.D. homogenea}
\end{gather*}


E.D. Lineales 


\begin{gather*}
a_{n}(x)\frac{d^{n}y}{dx^{n}}+a_{n-1}(x)\frac{d^{n-1}y}{dx^{n-1}}+......a_{1}(x)\frac{dy}{dx}+ a_{0}(x)y = g(x)
\end{gather*}



\begin{gather*}
n=1 \\
a_{1}(x)\frac{dy}{dx}+a_{0}(x)y=g(x) \\
\text{El objetivo, va  ser dejar suelta el termino de la diferencial} \\
\text{Se puede dividir entre a}_{1}(x) \\
\text{\colorbox[RGB]{248,231,28}{$\frac{a_{1}(x)}{a_{1}(x)}$}}\frac{dy}{dx}+\frac{a_{0}(x)}{a_{1}(x)}y=\frac{g(x)}{a_{1}(x)} \\
\vspace{0.5em} \\
P(x) =\frac{a_{0}(x)}{a_{1}(x)}  ;    Q(x) = \frac{g(x)}{a_{1}(x)} \\
\frac{dy}{dx}+P(x)y=Q(x)
\end{gather*}


E.D. Bernoulli $\rightarrow$  es una variacion a la forma Lineal


\begin{gather*}
\frac{dy}{dx}+P(x)y=Q(x)y^{n} \\
n<> 0,1
\end{gather*}


E.D. Exactas


\begin{gather*}
M(x,y)dx +N(x,y)dy = 0 \\
\text{es Exacta cuando se cumple:} \\
\frac{\partial M(x,y)}{\partial y} =\frac{\partial N(x,y)}{\partial x} 
\end{gather*}





\end{document}
